% LaTeX companion of 06-normal-subgroups-and-isomorphisms.typ. Both formats are editable.
\chapter{Normal subgroups and isomorphisms}\label{normal-subgroups-and-isomorphisms}

This chapter is a source-language transcription of `WorkSheets/412-WS24-Mywork.pdf', pp. 1--3. The page divisions below are part of the provenance: no theorem statement or proof cue is supplied from the reference-only PDFs.

\section{Kernels, quotients, and the first isomorphism theorem}

\textbf{Source transcription --- WS24, p. 1, Thm 8.16.} If \(f:G\rightarrow H\) is a group hom, then \(\ker f\) is a subgroup of \(G\). The handwritten proof first observes that for \(a,b \in \ker f\), \(f(a) = e_{H} = f(b)\) and hence \(f\left( {ab} \right) = f(a)f(b) = e_{H}\), so \(ab \in \ker f\). It then checks subgroup closure in the form: for \(g \in G\) and \(k \in \ker f\),

\[f\left( {g^{- 1}kg} \right) = {f(g)}^{- 1}f(k)f(g) = {f(g)}^{- 1}e_{H}f(g) = e_{H},\]

so \(g^{- 1}kg \in \ker f\). (The original Chinese line says 首先，group hom 的 ker 一定是 subgroup of \(G\)'\,` and then 然后我们证明 \(\ker f \lhd G\)'\,`.)

\textbf{Source transcription --- WS24, p. 1, Thm 8.17 and 8.18.}

\(\ker f = \left\{ e_{G} \right\}\) iff \(f\) is injective.

The sheet marks this as 已证过千遍.`\,' If \(N \lhd G\), then \(\pi:G\rightarrow\frac{G}{N}\) is a surjective group hom and \(\ker\pi = N\). It explicitly checks \(\pi\left( {g_{1}g_{2}} \right) = g_{1}g_{2}N = \left( {g_{1}N} \right)\left( {g_{2}N} \right)\), writes that every coset is \(Na\) for some \(a \in G\) and therefore is hit by \(\pi\), and notes \(\pi(a) = Ne = N\) iff \(a \in N\).

\textbf{Source transcription --- WS24, p. 1, Lemma 8.19.} For a group hom \(f:G\rightarrow H\) with \(\ker f = K\),

\[f(a) = f(b)\ \text{if and only if}\ Ka = Kb.\]

The source's forward implication is \(f\left( {ab^{- 1}} \right) = e_{H}\), hence \(ab^{- 1} \in K\) and \(a \equiv b\left( {\operatorname{mod}K} \right)\); for the converse, \(Ka = Kb\) gives \(ab^{- 1} \in K\), then \(f\left( {ab^{- 1}} \right) = e_{H}\) and, using \(f(a) = f(b)\) (the page annotates the equivalence with the reverse-multiplying calculation).

\begin{theorem}{\kn{First isomorphism theorem}}

If \(f:G\rightarrow H\) is a surjective group homomorphism, then

\[\frac{G}{\ker}f \sim = H.\]

\end{theorem}

\begin{proof}

\textbf{Source transcription --- WS24, p. 1, Thm 8.20.} Consider \(\psi:\frac{G}{\ker}f\rightarrow H\), sending \(Ka\mapsto f(a)\). It is well-defined because \(Ka = Kb\Rightarrow ab^{- 1} \in K\Rightarrow f(a) = f(b)\). It is injective by the preceding lemma, and it is surjective because every \(x \in H\) is \(f(g)\) for some \(g \in G\) when \(f\) is surjective. The source calls this 第一同构定理'\,` and annotates the displayed conclusion with the exceptional hypothesis that \(f\) must be surjective.

\end{proof}

\textbf{Source transcription --- WS24, p. 1, Thm 8.21.} If \(N \lhd G\), \(K\) is a subgroup of \(G\), and \(N \subset K\), then \(\frac{K}{N}\) is a subgroup of \(\frac{G}{N}\). The proof starts with \(gN \in \frac{G}{N}\) and \(kN \in \frac{K}{N}\); normality gives \(g^{- 1}kg \in K\), hence \(\left( {Ng} \right)^{({- 1}\}}\left( {Nk} \right)\left( {Ng} \right)\) (as written on the page) lies in \(\frac{K}{N}\). The source adds the Chinese reminder: 而如果 \(K \lhd G\)，则结论 更强: \(\frac{K}{N} \lhd \frac{G}{N}\)；但如果结论是包含 \(G\) 本身，和第七条类似.`\,'

\section{Second and third isomorphism theorems}

\textbf{Source transcription --- WS24, p. 2, Thm 8.22 (Third Isomorphism Theorem).} If \(N \lhd G\), \(K \lhd G\), and \(N \subset K\), then

\[\frac{K}{N} \lhd \frac{G}{N}\quad \land \quad\frac{\frac{G}{N}}{\frac{K}{N}} \sim = \frac{G}{K}.\]

The source begins the normality check with \(\left( {Ng} \right)^{({- 1}\}}\left( {Nk} \right)\left( {Ng} \right) = N\left( {g^{- 1}kg} \right)\) and, because \(K\) is normal in \(G\), \(g^{({- 1}\}}kg \in K\). For the quotient isomorphism it considers \(\pi:\frac{G}{N}\rightarrow\frac{G}{K}\), \(Na\mapsto Ka\), calling it an easy group hom and surjective. The ker note says: 即所有 \(a \in K\) 中等类的 \(N\)-cosets'\,`, so \(\ker\pi = \frac{K}{N}\), and the first isomorphism theorem yields the result.

\textbf{Source transcription --- WS24, p. 2, Second Isomorphism Theorem (group), Diamond Thm'\,`.} Let \(G\) be a group, \(S\) a subgroup of \(G\), and \(N \lhd G\). Then:

\begin{enumerate}
\tightlist
\item
  \(SN\) is a subgroup of \(G\);
\item
  \(N \lhd SN\);
\item
  \(S \cap N \lhd S\); and
\item
  \(\frac{SN}{N} \sim = \frac{S}{S \cap N}\).
\end{enumerate}

The source draws the diamond \(G\) over \(SN\), with \(S\) and \(N\) below and \(S \cap N\) at the base. It defines \(\psi:S\rightarrow S\frac{N}{N}\) by \(s\mapsto sN\); its kernel is \(\left\{ {s \in S\ \text{and}\ s \in N} \right\} = S \cap N\), so the first isomorphism theorem proves \(\frac{S}{S \cap N} \sim = S\frac{N}{N}\).

\textbf{Source transcription --- WS24, p. 2, Fourth Isomorphism Theorem (group), Lattice Thm'\,`.} With \(N \lhd G\), let \(\mathcal{G}\) be all subgroups of \(G\) containing \(N\) and \(\mathcal{N}\) all subgroups of \(\frac{G}{N}\). The source states \(\mathcal{G} \sim = \mathcal{N}\) by \(A\mapsto\frac{A}{N}\) and gives the correspondence cues 所有 \(\frac{G}{N}\) 的 subgroup \(T = \left\{ \frac{H}{N} \right\}\) for some \(H < G\)'\,` and, for a subgroup \(T < \frac{G}{N}\), choose \(H = \left\{ {a \in G:Na \in T} \right\}\), then prove \(H < G\) and \(\frac{H}{N} = T\).

\textbf{Source transcription --- WS24, p. 2, ring analogues.} 类比地 ring 也有四个 isomorphic thms.`\,' The page records the First Isomorphism Theorem for rings: if \(\varphi:R\rightarrow S\) is a ring hom, then \(\ker\varphi\) is a subring and an ideal, \(\operatorname{im}\varphi\) is a subring, and \(\operatorname{im}\varphi \sim = \frac{R}{\ker}\varphi\) (the page annotates the surjective case 虽然不说，但如果 \(\varphi\) surj，那么 \(\frac{R}{\ker}\varphi \sim = S\)'\,`). The Second Isomorphism Theorem for rings: if \(S\) is a subring of \(R\) and \(I\) an ideal of \(R\), then \(S + I = \left\{ {s + i:s \in S,i \in I} \right\}\) is a subring, \(S \cap I\) is an ideal of \(R\), and \(\frac{S + I}{I} \sim = \frac{S}{S \cap I}\).

\section{Fourth isomorphism theorem, simple groups, and finite abelian groups}

\textbf{Source transcription --- WS24, p. 3, Third and Fourth Isomorphism Theorems (ring).} If \(R\) is a ring and \(I\) an ideal of \(R\), the source lists:

\begin{enumerate}
\tightlist
\item
  for a subring \(A\) of \(R\), \(A + I\) is a subring of \(R\);
\item
  every subring of \(\frac{R}{I}\) is \(\frac{A}{I}\) for a subring \(A\) of \(R\);
\item
  if \(J\) is an ideal of \(R\) containing \(I\), then \(\frac{J}{I}\) is an ideal of \(\frac{R}{I}\);
\item
  every ideal of \(\frac{R}{I}\) is \(\frac{J}{I}\) for an ideal \(J\) of \(R\); and
\item
  \(\frac{R}{I}\) is isomorphic to \(\frac{R}{J}\) when \(\frac{J}{I}\) is the intervening ideal.
\end{enumerate}

The page's Fourth Isomorphism Theorem for rings is phrased: if \(I\) is an ideal of \(R\), define \(\mathcal{G}\) as all subrings of \(R\) containing \(I\) and \(\mathcal{N}\) as all subrings of \(\frac{R}{I}\); then \(\mathcal{G} \sim = \mathcal{N}\) under \(A\mapsto\frac{A}{I}\).

\textbf{Source transcription --- WS24, p. 3.} Corollary 8.23 says: if \(N\) is normal in \(G\), \(K\) is a subgroup of \(G\), and \(K\) contains \(N\), then \(K \lhd G\) iff \(\frac{K}{N} \lhd \frac{G}{N}\). The proof uses the Third Isomorphism Theorem in one direction and, in the other, for \(g \in G\), \(k \in K\), writes \(\left( {Ng} \right)^{({- 1}\}}\left( {Nk} \right)\left( {Ng} \right) = Nk'\) for some \(k' \in K\), hence \(g^{({- 1}\}}kg = Nt\) for \(t \in K\); since \(N \subset K\), this lies in \(K\).

The definition is retained verbatim in meaning: A group \(G\) is simple iff 它有且只有 \(\left\{ e_{G} \right\}\) 和 \(G\) 自己这两个 normal subgroup.`\,' The sheet states \(G\) 为 simple abelian group iff \(G \sim = \mathbb{Z}_{p}\) for some prime \(p\).`\,' It finishes with the Fundamental Structure Theorem for finite Abelian groups:

\[G \sim = \mathbb{Z}_{p_{1}^{a_{1}}} \times \mathbb{Z}_{p_{2}^{a_{2}}} \times \ldots \times \mathbb{Z}_{p_{n}^{a_{n}}},\]

where \(p_{1},\ldots,p_{n}\) are prime numbers (可以重复'\,`), and the isomorphism is unique up to reordering.
